% !TEX program = lualatex
% !TeX encoding = UTF-8
\PassOptionsToPackage{unicode}{hyperref}
\PassOptionsToPackage{bookmarksnumbered,bookmarksopen}{hyperref}
\documentclass[12pt,a4paper]{article}

% ============================================================
% 日本語設定 (LuaLaTeX + luatexja)
% ============================================================
\usepackage{luatexja}
\usepackage{luatexja-fontspec}

% ラテン文字フォント
\setmainfont{Times New Roman}[Ligatures=TeX]
\setsansfont{Arial}[Ligatures=TeX]
\setmonofont{Consolas}[Scale=MatchLowercase]

% 日本語フォント – Yu Gothic (Windows に標準搭載)
\setmainjfont{Yu Gothic}[UprightFont = * Regular, BoldFont = * Bold]
\setsansjfont{Yu Gothic}[UprightFont = * Regular, BoldFont = * Bold]
\setmonojfont{Yu Gothic}[UprightFont = * Regular, BoldFont = * Bold]

% ============================================================
% その他のパッケージ
% ============================================================
\usepackage{amsmath,amssymb,amsthm}
\usepackage{graphicx}
\usepackage{hyperref}
\usepackage{url}
\usepackage{xcolor}
\usepackage{booktabs}
\usepackage{geometry}
\geometry{margin=2.5cm}

% ============================================================
% YUCT マクロ
% ============================================================
\newcommand{\Keff}{K_{\mathrm{eff}}}
\newcommand{\kappac}{\kappa_c}
\newcommand{\eps}{\varepsilon}
\newcommand{\betagamma}{\beta\gamma}
\newcommand{\betac}{\beta}
\newcommand{\tauf}{\tau}

% ============================================================
% 定理環境（日本語）
% ============================================================
\newtheorem{theorem}{定理}
\newtheorem{lemma}{補題}
\newtheorem{corollary}{系}
\newtheorem{definition}{定義}
\newtheorem{prediction}{予測}

% ============================================================
% ハイパーリンク設定
% ============================================================
\hypersetup{
	colorlinks=true,
	linkcolor=blue,
	citecolor=blue,
	urlcolor=blue,
	pdftitle={付録 S: 冷却原子における光子の負の時間遅延 — 協調量子ダイナミクスの発現として：Yakushev 統一調整理論 (YUCT) に基づく改訂解析},
	pdfauthor={Alexey V. Yakushev},
	pdfkeywords={YUCT, 調整効率, 負の群遅延, 光子–原子相互作用, 対数スケーリング, 温度独立性, 撤回}
}

\begin{document}
	
	\begin{titlepage}
		\centering
		\vspace*{1cm}
		{\Huge\bfseries 付録 S: 冷却原子における光子の負の時間遅延 — 協調量子ダイナミクスの発現として：\\ Yakushev 統一調整理論 (YUCT) に基づく改訂解析}
		\vspace{1.5cm}
		
		{\large Alexey V. Yakushev\textsuperscript{1}}
		\vspace{0.3cm}
		{\normalsize \textsuperscript{1}Yakushev Research, \\ 
			\textbf YUCT 理論: \url{https://yuct.org/}, \\ 
			\textbf YPSDC プロトコル: \url{https://ypsdc.com/}\\
			\textbf DOI: \url{https://doi.org/10.5281/zenodo.18444598}}
		\vspace{1cm}
		{\normalsize 2026年5月 \\ \large バージョン V.3.0 \\(修正された対数スケーリングと13\%温度予測の撤回)}
		
		\vfill
		
		\begin{abstract}
			\noindent
			最近の画期的な実験 (Angulo 他, 2024) は、超冷却原子雲を通過する光子が負の時間遅延で媒質を出ることができることを示した – 波束のピークが入射したよりも早く現れる。この直観に反する現象は因果律を侵害しないものの、深い量子干渉効果を明らかにする。本研究では、Yakushev 統一調整理論 (YUCT) に基づく説明を再検討する。普遍スケーリング則 \(\eps = \kappac \alpha (\ln\Keff)^{\betac}\) (\(\betac = 2/3\), \(\kappac = 1/3\)，付録 X で導出) に基づき、群遅延 \(\tau_g\) と調整効率 \(\Keff\) の正しい関係が対数的であることを示す: \(|\tau_g| = \alpha_0 \ln\Keff\)。負の時間遅延の存在は既に実験によって確認されており; \((-1.14 \pm 0.22)\tau_0\) もの値が観測されている。しかし、以前の線形仮定とは対照的に、対数形式は \(\tau_g\) の温度依存性が極めて弱いことを意味する。標準的な温度劣化関係 \(\Keff(T) = K_0 (T/T_0)^{-\gamma}\) (\(\gamma \approx 0.30\), 地球物理学的データから導出) を用いると、温度を2倍にしたときの \(|\tau_g|\) の相対変化が \(10^{-8}\) 程度であることが分かり、これは現在の実験の感度をはるかに下回る。従って、以前に主張された13\%効果は誤った線形モデルのアーティファクトであり、撤回されなければならない。我々は修正された理論的枠組みを提供し、光学的深さやパルス持続時間の変化など、対数スケーリングに敏感であり得る代替実験的検証を議論する。YUCT の核心的予測 – 負の群遅延が調整効率の発現であること – は依然として有効であるが、冷却原子系におけるその検証可能性は改良されたセットアップを必要とする。
		\end{abstract}
		
		\vspace{0.5cm}
		\noindent\textbf{キーワード:} YUCT, 調整効率, 負の群遅延, 光子–原子相互作用, 対数スケーリング, 温度独立性, 撤回.
		
	\end{titlepage}
	
	\tableofcontents
	\newpage
	
	\section{はじめに}
	
	共鳴媒質中の光の伝播は、数十年にわたって量子光学の要となってきた。群遅延、パルス整形、超光速効果は理論的および実験的によく理解されている \cite{Brillouin1960, Milonni2004}。それにもかかわらず、根本的な疑問が残る：光子が原子励起として過ごす時間の物理的意味は何か？
	
	Angulo 他 \cite{Angulo2024} による画期的な実験は、透過光子が別のプローブビームに誘起する交差カー位相シフトを測定することでこの問題に取り組んだ。彼らは、透過光子に対する積分原子励起時間 \(\tau_T\) が負になり得ること、そして \(\tau_g < 0\) の場合でも群遅延 \(\tau_g\) と一致することを発見した。具体的には、実験は \(\tau_T/\tau_0 = -1.14 \pm 0.22\) (10 ns パルス, OD~4) および \(0.54 \pm 0.28\) (異なるパラメータ) の値を報告した \cite{Angulo2024, Nixon2024}。この驚くべき結果は単純な直観に挑戦し、より深い理論的枠組みを要請する。
	
	Yakushev 統一調整理論 (YUCT) \cite{Yakushev2026a} はそのような枠組みを提供する。YUCT は、量子光学から重力に至るまで全ての物理的相互作用が、\textbf{調整効率} \(\Keff\) によって定量化される基本的な調整プロセスの発現であると仮定する。普遍スケーリング則は任意の相対誤差 \(\eps\) を \(\Keff\) と結びつける (付録 X, 式~\ref{eq:error-law-corrected} 参照):
	\begin{equation}
		\eps = \kappac \alpha (\ln\Keff)^{\betac}, \qquad \betac = \frac{2}{3},\ \kappac = \frac{1}{3},
		\label{eq:scaling}
	\end{equation}
	ここで \(\alpha\) はシステム依存である。この法則は DNA 複製誤差から宇宙マイクロ波背景放射の変動に至るまで、40 桁以上にわたって検証されている \cite{Yakushev2026c}。光–物質相互作用の文脈では、\(\Keff\) は光子と原子集団間のコヒーレンスを測定する。
	
	ここで我々は、負の時間遅延実験が YUCT の直接的な結果であることを示す。しかし、\(\tau_g\) と \(\Keff\) の間に線形関係を仮定した以前の試みを修正しなければならない。YUCT の情報理論的構造から導出される正しい関係は対数的である。これにより、予測される温度依存性が劇的に変化し、現在の実験精度では無視できるほどになる。従って、我々は以前に発表された13\%予測を撤回し、修正された自己無撞着な解析を提供する。
	
	\section{YUCT の基本概念}
	
	\subsection{調整効率 \(\Keff\)}
	
	YUCT では、任意の相互作用系は辞書 \(D\) (共有プロトコル、状態) と、特定の調整動作を活性化するインデックス \(i\) によって特徴づけられる。調整の効率は
	\begin{equation}
		\Keff = \frac{H(\text{動作})}{H(\text{インデックス})},
	\end{equation}
	ここで \(H\) はシャノンエントロピーを表す。量子系では \(\Keff\) は非常に大きくなり得、完全なエンタングルメントの極限で無限大に近づく。
	
	\subsection{普遍スケーリング則}
	
	変動と誤差は式~\eqref{eq:scaling} に従う。この法則は辞書のフラクタル幾何から生じ、中性子崩壊から銀河回転曲線に至るまで多様な系で経験的に確認されている \cite{Yakushev2026c}。対数形式は必須である：これはインデックス空間における有効分解能が辞書サイズに対して対数的にしか成長しないという事実を反映している。
	
	\subsection{\(\Keff\) の温度依存性}
	
	多くの物理系では、熱的変動のため調整効率は温度とともに減少する。臨界現象とフラクタルスケーリングから、我々は
	\begin{equation}
		\Keff(T) = K_0 \left(\frac{T}{T_0}\right)^{-\gamma}, \quad \gamma > 0,
		\label{eq:K_T}
	\end{equation}
	を持ち、ここで \(\gamma\) は材料依存の指数である。積 \(\betac\gamma\) は地球物理学的データから独立に決定されている: \(\betagamma = 0.20 \pm 0.03\) \cite{Yakushev2026b}。これは \(\gamma = \betagamma / \betac = 0.20 / (2/3) = 0.30\) を意味する。この値は非秩序系に対する一般的なスケーリング議論と整合する。
	
	\section{光子–原子相互作用の YUCT モデル}
	
	冷却原子集団に入射する共鳴光子パルスを考える。原子は \textbf{辞書} \(D\) として機能する：それらは特定の量子状態に準備される (例えばレーザー冷却とトラッピングによって)。光子は \textbf{インデックス} – その波束形状と周波数 – を運ぶ。相互作用は、光子が透過、吸収、または散乱され得る調整状態をもたらす。
	
	Angulo 他 \cite{Angulo2024} の実験で興味対象は、透過光子に対する \textbf{原子励起時間} \(\tau_T\) である。YUCT から、この時間は群遅延 \(\tau_g\) に比例する。Thompson 他 \cite{Thompson2023} の弱値解析は、関連領域において透過光子に対して \(\tau_T = \tau_g\) であることを示している。
	
	決定的なステップは \(\tau_g\) を \(\Keff\) と結びつけることである。群遅延は光子と原子集団の間のコヒーレントな相互作用から生じる。この相互作用の大きさは \(\Keff\) 自体 (巨大になり得る) に比例するのではなく、光子が原子状態について運ぶ \textbf{情報} の量に比例する。YUCT では、関連する情報量は調整効率の対数である。なぜならインデックス長 \(\ell = \log_2 M\) が識別可能な状態の数を決定するからである。従って我々は仮定する：
	\begin{equation}
		|\tau_g| = \alpha_0 \, \ln\Keff,
		\label{eq:tau_log}
	\end{equation}
	ここで \(\alpha_0\) は時間の次元を持つ基本時間定数である。ルビジウム D2 線に対して、自然線幅 \(\Gamma = 2\pi \times 6\) MHz は特性時間 \(\tau_{\text{nat}} = 1/\Gamma \approx 26\) ns を与える。\(K_0 \sim 10^4\) での実験値 \(|\tau_g| \approx 30\) ns にフィッティングすると \(\alpha_0 = 30\,\text{ns} / \ln(10^4) = 30 / 9.21 \approx 3.26\) ns を得る。これは妥当である：素相互作用時間は数ナノ秒程度である。
	
	\subsection{群遅延の温度依存性}
	
	式~\eqref{eq:K_T} を式~\eqref{eq:tau_log} に代入すると
	\begin{equation}
		|\tau_g(T)| = \alpha_0 \ln\!\left(K_0 \left(\frac{T}{T_0}\right)^{-\gamma}\right)
		= \alpha_0 \ln K_0 - \alpha_0 \gamma \ln\!\left(\frac{T}{T_0}\right).
	\end{equation}
	\(|\tau_g(T_0)| = \alpha_0 \ln K_0\) と書けば、我々は
	\begin{equation}
		|\tau_g(T)| = |\tau_g(T_0)| - \alpha_0 \gamma \ln\!\left(\frac{T}{T_0}\right).
		\label{eq:tau_T_log}
	\end{equation}
	を得る。これが中心的な結果である：群遅延の大きさは温度の増加とともに \textbf{対数的} に減少する。温度比 \(T/T_0 = 2\) に対する変化は
	\begin{equation}
		\Delta |\tau_g| = \alpha_0 \gamma \ln 2.
	\end{equation}
	\(\alpha_0 \approx 3.26\) ns と \(\gamma = 0.30\) を用いると、\(\Delta |\tau_g| \approx 3.26 \times 0.30 \times 0.693 \approx 0.68\) ns を得る。基準値 \(|\tau_g(T_0)| \approx 30\) ns に対する相対変化は \(0.68 / 30 \approx 2.3\%\) である。これは依然として実験の統計的不確実性 (\(\pm 0.22\) 単位 \(\tau_0\)，すなわち約 \(\pm 5.7\) ns) の範囲内にある。しかし、この効果は検出可能性の限界にある。より重要なことは、変化が \textbf{対数的} であり、べき乗則ではなく、その大きさが以前に主張された13\%よりはるかに小さいことである。
	
	\subsection{以前の (誤った) 線形モデルとの比較}
	
	以前の研究では、線形関係 \(|\tau_g| \propto \Keff\) が仮定されていた。これは \(|\tau_g| \propto T^{-\gamma}\) をもたらし、\(\gamma = 0.30\) で温度を2倍にすると13\%の変化をもたらした。この予測は量子系における \(\Keff\) の役割の誤解に基づいていた。正しい対数関係は、\(\Keff\) がエントロピーの比であり、直接的な物理的場の振幅ではないために生じる。線形モデルは YUCT の情報理論的基礎と両立せず、ここに撤回される。
	
	\section{既存データからの数値的整合性チェック}
	
	Angulo 他 \cite{Angulo2024} の実験は、\(\alpha_0\) と \(K_0\) を推定することを可能にする特定のデータ点を提供する。10 ns パルス、光学深さ OD~4 の構成に対して、報告された原子励起時間は \(\tau_T = -1.14 \tau_0\) であった。\(\tau_0 \approx 26\) ns を取れば、\(|\tau_g| \approx 30\) ns を得る。\(\Keff(T_0) = K_0\) が \(10^4\) 程度であると仮定すれば (典型的な冷却原子集団は \(10^6\) 個の原子と高いコヒーレンスを持つ)，\(\alpha_0 = 30\,\text{ns} / \ln(10^4) \approx 3.26\) ns を得る。この値は基本的な光子–原子相互作用の特性時間尺度として合理的である。追加の自由パラメータは必要ない。
	
	\section{標準量子光学との比較}
	
	標準量子光学の形式では、共鳴パルスの群遅延 \(\tau_g\) は透過位相の周波数微分 \(\tau_g = d\phi/d\omega\) で与えられる。ローレンツ原子応答に対して、これは異常分散のために共鳴の翼で負になり得る特性形状を与える \cite{Milonni2004}。しかし、標準理論はドップラー広がり (温度で線形、対数ではない) などの自明な効果を除いて、\(\tau_g\) の原子雲温度への依存性を予測しない。観測される負の時間遅延はエネルギー交換を含まないコヒーレント干渉効果として理解される。
	
	YUCT は遅延の大きさを調整効率 \(\Keff\) と結びつけることで標準理論を超える。対数依存性は特徴的なシグネチャである：原子密度、光学深さ、またはパルス持続時間を変化させることで \(\Keff\) を変えられれば、遅延は \(\ln\Keff\) としてスケールするはずである。これは温度変化に依存しない検証可能な予測である。
	
	\section{修正された実験的提案}
	
	温度依存性が極めて弱いため (対数的、温度2倍で相対変化が数パーセント程度)、現在の精度で検出される可能性は低い。そこで我々は、より制御された方法で \(\Keff\) を直接変化させる代替実験を提案する。
	
	\subsection{光学深さの変化}
	
	調整効率 \(\Keff\) は集団内の原子数と相互作用強度に比例する。光学深さを減らすことで (例えば原子密度を下げる)，\(\Keff\) は減少する。群遅延の大きさは \(|\tau_g| = \alpha_0 \ln\Keff\) に従うはずである。\(\Keff\) を10分の1に減らすと (例えば \(10^4\) から \(10^3\) へ)，\(|\tau_g|\) の変化は \(\alpha_0 (\ln 10^4 - \ln 10^3) = \alpha_0 \ln 10 \approx 3.26 \times 2.3026 \approx 7.5\) ns となり，基準値の約25\%で容易に測定可能である。このような実験は単純である：異なる原子密度に対して \(\tau_T\) の測定を繰り返し (他の全てのパラメータは固定)，対数スケーリングを検証する。
	
	\subsection{パルス持続時間の変化}
	
	光子パルスの情報内容 (インデックス長) はその時間的持続時間と帯域幅に依存する。より短いパルスはより多くのスペクトル情報を運ぶため、より大きな \(\Keff\) を持つ。パルス持続時間を 1 ns から 20 ns まで系統的に変化させることで、群遅延が \(\ln(\text{帯域幅})\) に従ってどのようにスケールするかを観察できる。これは対数則のもう一つの直接的な検証を提供する。
	
	\subsection{温度依存性をヌルテストとして}
	
	予測される温度効果は小さいが、正確にゼロではない。ヌル測定 (実験誤差内で検出可能な変化がない) は我々の対数モデルと整合し、以前の線形モデルは明確な13\%変化を予測したであろう。従って、既存のデータは既に線形モデルを支持しておらず、将来の高精度実験は最終的に対数シフトを検出するかもしれない。我々は実験家に両方を検証することを奨励する。
	
	\section{光子質量予測との関連：整合性チェック}
	\label{sec:photon_mass_consistency}
	
	YUCT フレームワークはまた、同一の質量梯子 \(m = m_e q^{N_f}\) (\(q=(3/2)^{1/3}\)) から導出された光子の静止質量に対する盲目的予測を生み出した。実験的上限と整合する最適ノードは \(N_\gamma = -410\) であり、これは
	\[
	m_\gamma \approx 7.3 \times 10^{-19}~\text{eV} \quad (\text{付録 Photon Mass 参照})
	\]
	を与える。
	
	このような微小な非ゼロ光子質量が冷却原子中の群遅延測定に影響を与えるかもしれないと疑問に思うかもしれない。\(m_\gamma\) に対応するコンプトン波長は
	\[
	\lambda_c = \frac{\hbar}{m_\gamma c} \approx 2.7\times10^{11}~\text{m},
	\]
	であり、これは原子雲 (典型的なサイズ \(\sim 1\) mm) よりもはるかに大きい。従って、有限光子質量による伝播効果はこの実験では完全に無視できる。負の時間遅延は純粋な調整効果であり、絶対質量スケールとは独立である。
	
	逆に、単一のフレームワークが以下の両方を正しく予測するという事実は：
	\begin{itemize}
		\item \(\ln\Keff\) によって大きさが設定される負の群遅延の存在 (本付録)，
		\item 光子質量の上限 (\(m_\gamma < 10^{-18}\) eV) とその離散ノード \(N_\gamma = -410\) (付録 Photon Mass)，
	\end{itemize}
	強力な整合性チェックを構成する。二つの予測は非常に異なる物理領域 (量子光学 vs. 基本粒子特性) に関わるが、同じ基本定数 \(q\), \(S_{\mathrm{odd}}/S_{\mathrm{even}}\)，そして調整効率 \(\Keff\) の概念を共有する。それらを整合させるために特別な調整は必要ない。
	
	従って、負の光子遅延実験は YUCT 光子質量予測と矛盾せず；むしろ両者は統合調整パラダイムを支持する独立した柱として機能する。
	
	\section{以前の盲目的予測の撤回}
	
	本付録の以前のバージョン (V.2.0) で、我々は群遅延が \(T^{-0.20}\) としてスケールし、温度2倍で13\%の変化をもたらすと誤って予測した。この予測は誤った線形関係 \(|\tau_g| \propto \Keff\) に基づいていた。ここで導出された正しい対数関係は、温度効果が約2桁小さいことを示している。従って、13\%効果の主張は無効であり、正式に撤回される。この見落としについてお詫びし、不一致を特定してくれた数学的パッチの共著者に感謝する。
	
	それにもかかわらず、YUCT の核心的概念 – 負の群遅延が調整効率の発現であること – は依然として有効である。負の原子励起時間の実験的確認は、光子と原子が調整された辞書を共有するという考えを既に支持している。修正されたスケーリング則は今や将来の検証のための洗練された理論的枠組みを提供する。
	
	\section{結論}
	
	我々は負の光子遅延実験に対する YUCT 解析を修正した。YUCT の普遍誤差則は対数的であり、従って群遅延と調整効率の関係は対数的でなければならない: \(|\tau_g| = \alpha_0 \ln\Keff\)。これは非常に弱い温度依存性 (対数的、温度2倍で相対変化が数パーセント程度) をもたらし、以前に主張された13\%効果を除去する。誤った予測は撤回される。光学深さやパルス持続時間の変化などの代替検証は、対数スケーリングに対してはるかに敏感なプローブを提供する。負の時間遅延の存在は量子相互作用の調整された性質の驚くべき確認として残り、YUCT は一貫した情報理論的解釈を提供する。
	
	\section*{謝辞}
	
	著者は、線形仮定と YUCT の対数構造との間の不一致を指摘してくれた数学的パッチの共著者に感謝する。また、YUCT セミナーの参加者に刺激的な議論を、そしてトロントグループ (Aephraim Steinberg, Daniela Angulo とその同僚) に彼らの美しい実験に感謝する。
	
	\section*{データ利用可能性}
	
	全ての計算および追加資料は \url{https://github.com/Alexey-Yakushev-YUCT/YPSDC} で入手可能である。
	
	\begin{thebibliography}{99}
		
		\bibitem{Angulo2024}
		D. Angulo, K. Thompson, A. Jiao, H. Wiseman, A. Steinberg, ``Experimental evidence that a photon can spend a negative amount of time in an atom cloud,'' arXiv:2409.03680 [quant-ph] (2024).
		
		\bibitem{Nixon2024}
		V.-M. Nixon et al., ``Measuring the time transmitted photons spend as atomic excitations,'' American Physical Society DAMOP Conference, Session Y07 (June 2024).
		
		\bibitem{Yakushev2026a}
		A. V. Yakushev, ``Yakushev Unified Coordination Theory (YUCT),'' Zenodo, \url{https://doi.org/10.5281/zenodo.18444599} (2026).
		
		\bibitem{Yakushev2026b}
		A. V. Yakushev, ``Appendix P: Temperature Dependence of Gravity,'' Zenodo (2026).
		
		\bibitem{Yakushev2026c}
		A. V. Yakushev, ``Appendix L: Fractal Error Scaling,'' Zenodo (2026).
		
		\bibitem{Yakushev2026d}
		A. V. Yakushev, ``Appendix X: Generalization of Shannon Information Theory,'' Zenodo (2026).
		
		\bibitem{Thompson2023}
		K. Thompson et al., ``How much time does a photon spend as an atomic excitation before being transmitted?'' arXiv:2310.00432 (2023).
		
		\bibitem{Brillouin1960}
		L. Brillouin, \textit{Wave Propagation and Group Velocity} (Academic Press, 1960).
		
		\bibitem{Milonni2004}
		P. Milonni, \textit{Fast Light, Slow Light and Left-Handed Light} (CRC Press, 2004).
		
		\bibitem{Wiseman2023}
		H. M. Wiseman, A. M. Steinberg, and M. Hallaji, ``Obtaining a single-photon weak value from experiments using a strong coherent state,'' AVS Quantum Sci. \textbf{5}, 024401 (2023).
		
	\end{thebibliography}
	
\end{document}